\documentclass{ximera}

\input{../../preamble.tex}

\title{Double Angles}

\begin{document}

\begin{abstract}
changing angles
\end{abstract}
\maketitle



The sum rules for sine and cosine give us a way to talk about double angles.




\[ \cos(\theta_1 + \theta_2) =  \cos(\theta_1) \cdot \cos(\theta_2) - \sin(\theta_1) \cdot \sin(\theta_2) \]

and

\[ \sin(\theta_1 + \theta_2) =  \cos(\theta_1) \cdot \sin(\theta_2) + \sin(\theta_1) \cdot \cos(\theta_2) \]





If we let $\theta_1 = \theta_2$, then we get the ``\textbf{double angle rules}''.



\[ \cos(2 \, \theta) =  (\cos(\theta))^2 - (\sin(\theta))^2  \]

and

\[ \sin(2 \, \theta) =  2 \, \sin(\theta) \,  \cos(\theta) \]












\subsection*{More Relationships}


$\sin(x)$ and $\cos(x)$ are the coordinates for points on the unit circle.

Therefore, we know that 

\[ (\cos(x))^2 + (\sin(x))^2 = 1\]


From this equation, we also get 


\[ (\cos(x))^2 = 1 - (\sin(x))^2 \]

and

\[ 1 - (\cos(x))^2 = (\sin(x))^2 \]


Ssubstituting these into our equaiton for the double angle cosine rules gives two alternative.


\begin{align*}
\cos(2 \, \theta) & = (\cos(\theta))^2 - (\sin(\theta))^2 \\
& = (\cos(\theta))^2 - (1 - (\cos(x))^2)  \\
& = 2 \, (\cos(\theta))^2 - 1
\end{align*}

and


\begin{align*}
\cos(2 \, \theta) & = (\cos(\theta))^2 - (\sin(\theta))^2 \\
& = 1 - (\sin(\theta))^2 - (\sin(\theta))^2  \\
& = 1 - 2 \, (\sin(\theta))^2
\end{align*}


\textbf{\textcolor{blue!55!black}{Arithmetic}} 

Arithmetic with trigonometric function is kind of like arithmetic with fractions.  

Our arithmetic rules for fraction addition insist on common denominators.

With trigonometric functions, we need the arguments (the insides) to be equal.


The double angle rules give us a way convert angles to find common angles.

This is very helpful when solving equations.




\begin{example} Solving Equations


Solve $\cos(2x) = \sin(x)$

\textbf{\textcolor{red!75!green}{explanation}} 


The problem here is really the arguments of sine and cosine.  They are different.

We can make them the same using the double angle rules.


\[  \cos(2x) = \sin(x)  \]

\[  1 - 2(\sin(x))^2 = \sin(x)  \]

\[  0 = 2(\sin(x))^2 + \sin(x) - 1  \]

\[  0 = (2 \, \sin(x) - 1) (\sin(x) + 1)  \]


Either $0 = 2 \, \sin(x) - 1$ or $0 = \sin(x) + 1$.

Either $\frac{1}{2} = \sin(x) $ or $-1 = \sin(x)$.


On the interval $[0, 2\pi)$, we get three solutions.


\[ \frac{\pi}{6}, \frac{5\pi}{6}, \frac{3\pi}{2}  \]


These can be extended with a period of $2\pi$ to get all of the solutions.



\end{example}






















\begin{center}
\textbf{\textcolor{green!50!black}{ooooo-=-=-=-ooOoo-=-=-=-ooooo}} \\

more examples can be found by following this link\\ \link[More Examples of Trigonometric Forms]{https://ximera.osu.edu/csccmathematics/precalculus/precalculus/trigForms/examples/exampleList}

\end{center}






\end{document}
